% GENERATED by scripts/gen_appendix.py -- do not edit.

\subsection{The placement family (31 descriptors)}
\noindent The rigid-body pose and the two CDR3 loops, read in the groove frame.

\small
\begin{longtable}{@{}>{\raggedright\arraybackslash}p{0.19\textwidth}p{0.14\textwidth}p{0.12\textwidth}p{0.08\textwidth}p{0.34\textwidth}@{}}
\toprule
descriptor & class & operator & units & definition \\
\midrule\endfirsthead
\toprule descriptor & class & operator & units & definition \\\midrule\endhead
\texttt{pitch}$^{\dagger}$ & geometric & frame coordinate & $^\circ$ & Incident angle of the TCR out of the groove plane. **Banned as a feature**: it reproduces no clean geometric angle yet out-discriminates every one of them, which is AlphaFold-confidence contamination rather than geometry. \\
\texttt{crossing}$^{\dagger}$ & geometric & frame coordinate & $^\circ$ & Crossing (scanning) angle between the V$\alpha$->V$\beta$ axis projected into the groove plane and the groove long axis. \\
\texttt{crossing\_\allowbreak{}signed} & geometric & frame coordinate & $^\circ$ & The same angle on [-180, 180); its sign is the docking polarity, canonical or reversed. \\
\texttt{dock\_\allowbreak{}d} & geometric & frame coordinate & \AA & MHC-stub to TCR-stub rigid-body separation. \\
\texttt{dock\_\allowbreak{}torsion} & geometric & frame coordinate & rad & Rigid-body dihedral of the TCR about the MHC stub; the docking twist. Circular, wraps at +-pi. \\
\texttt{dock\_\allowbreak{}tcr\_\allowbreak{}uy} & geometric & frame coordinate & cosine & y component of the TCR stub unit vector in the MHC frame. \\
\texttt{dock\_\allowbreak{}tcr\_\allowbreak{}uz} & geometric & frame coordinate & cosine & z component of the TCR stub unit vector; how high the receptor body rides over the groove. \\
\texttt{dock\_\allowbreak{}mhc\_\allowbreak{}uy} & geometric & frame coordinate & cosine & y component of the MHC stub unit vector. \\
\texttt{dock\_\allowbreak{}mhc\_\allowbreak{}uz} & geometric & frame coordinate & cosine & z component of the MHC stub unit vector. \\
\texttt{height} & geometric & frame coordinate & \AA & Elevation of the CDR C$\alpha$ centroid above the groove plane. \\
\texttt{shift\_\allowbreak{}u} & geometric & frame coordinate & \AA & In-plane displacement of that centroid from the peptide centroid along the groove long axis. \\
\texttt{shift\_\allowbreak{}w} & geometric & frame coordinate & \AA & The same along the groove short axis. \\
\texttt{offset}$^{\dagger}$ & geometric & frame coordinate & \AA & Length of the in-plane displacement; lateral shift whatever its direction. \\
\texttt{cdr3a\_\allowbreak{}reach} & geometric & frame coordinate & \AA & Distance from the loop's C$\alpha$ centroid to the peptide C$\alpha$ centroid; how far CDR3$\alpha$ reaches. \\
\texttt{cdr3a\_\allowbreak{}ou} & geometric & frame coordinate & cosine & Where CDR3$\alpha$ sits over the groove, along the long axis u. \\
\texttt{cdr3a\_\allowbreak{}ow} & geometric & frame coordinate & cosine & Where CDR3$\alpha$ sits over the groove, along the short axis w. \\
\texttt{cdr3a\_\allowbreak{}on} & geometric & frame coordinate & cosine & Where CDR3$\alpha$ sits over the groove, along the groove normal n. \\
\texttt{cdr3a\_\allowbreak{}au} & geometric & frame coordinate & cosine & Orientation of CDR3$\alpha$'s N->C axis against the groove long axis u. \\
\texttt{cdr3a\_\allowbreak{}aw} & geometric & frame coordinate & cosine & Orientation of CDR3$\alpha$'s N->C axis against the groove short axis w. \\
\texttt{cdr3a\_\allowbreak{}an} & geometric & frame coordinate & cosine & Orientation of CDR3$\alpha$'s N->C axis against the groove normal n. \\
\texttt{cdr3a\_\allowbreak{}topep} & geometric & frame coordinate & \AA & Minimum C$\alpha$-C$\alpha$ distance from CDR3$\alpha$ to the peptide; its engagement depth. \\
\texttt{cdr3a\_\allowbreak{}ext} & geometric & frame coordinate & \AA & End-to-end extension of CDR3$\alpha$, the C$\alpha$\_N to C$\alpha$\_C distance. \\
\texttt{cdr3b\_\allowbreak{}reach} & geometric & frame coordinate & \AA & Distance from the loop's C$\alpha$ centroid to the peptide C$\alpha$ centroid; how far CDR3$\beta$ reaches. \\
\texttt{cdr3b\_\allowbreak{}ou} & geometric & frame coordinate & cosine & Where CDR3$\beta$ sits over the groove, along the long axis u. \\
\texttt{cdr3b\_\allowbreak{}ow} & geometric & frame coordinate & cosine & Where CDR3$\beta$ sits over the groove, along the short axis w. \\
\texttt{cdr3b\_\allowbreak{}on} & geometric & frame coordinate & cosine & Where CDR3$\beta$ sits over the groove, along the groove normal n. \\
\texttt{cdr3b\_\allowbreak{}au} & geometric & frame coordinate & cosine & Orientation of CDR3$\beta$'s N->C axis against the groove long axis u. \\
\texttt{cdr3b\_\allowbreak{}aw} & geometric & frame coordinate & cosine & Orientation of CDR3$\beta$'s N->C axis against the groove short axis w. \\
\texttt{cdr3b\_\allowbreak{}an} & geometric & frame coordinate & cosine & Orientation of CDR3$\beta$'s N->C axis against the groove normal n. \\
\texttt{cdr3b\_\allowbreak{}topep} & geometric & frame coordinate & \AA & Minimum C$\alpha$-C$\alpha$ distance from CDR3$\beta$ to the peptide; its engagement depth. \\
\texttt{cdr3b\_\allowbreak{}ext} & geometric & frame coordinate & \AA & End-to-end extension of CDR3$\beta$, the C$\alpha$\_N to C$\alpha$\_C distance. \\
\bottomrule
\end{longtable}

\subsection{The interface family (26 descriptors)}
\noindent How much contact there is, and of what chemical kind.

\small
\begin{longtable}{@{}>{\raggedright\arraybackslash}p{0.19\textwidth}p{0.14\textwidth}p{0.12\textwidth}p{0.08\textwidth}p{0.34\textwidth}@{}}
\toprule
descriptor & class & operator & units & definition \\
\midrule\endfirsthead
\toprule descriptor & class & operator & units & definition \\\midrule\endhead
\texttt{burial} & geometric & count or share & \AA$^2$ & Interface buried surface, SASA(TCR) + SASA(pMHC) - SASA(complex), by Shrake-Rupley. \\
\texttt{extent} & compositional & count or share & count & Distinct TCR residues contacting the pMHC over both receptor interfaces. \\
\texttt{chain\_\allowbreak{}balance} & compositional & count or share & fraction & min(a,b)/(a+b) over TCR:peptide contacts by chain; 0.5 when both chains engage equally, 0 when only one does. \\
\texttt{n\_\allowbreak{}contacts\_\allowbreak{}tp} & compositional & count or share & count & TCR-peptide residue-residue contacts. \\
\texttt{n\_\allowbreak{}contacts\_\allowbreak{}tm} & compositional & count or share & count & TCR-MHC residue-residue contacts. \\
\texttt{n\_\allowbreak{}pep\_\allowbreak{}contacted} & compositional & count or share & count & Distinct peptide residues the TCR contacts. \\
\texttt{n\_\allowbreak{}hbond} & compositional & count or share & count & Polar N/O atom pairs within 3.5 A across TCR:peptide. \\
\texttt{ct\_\allowbreak{}tp\_\allowbreak{}salt\_\allowbreak{}bridge}$^{\dagger}$ & compositional & count or share & count & Cationic-N / anionic-O pairs within 4 A across TCR:peptide. \\
\texttt{ct\_\allowbreak{}tp\_\allowbreak{}aromatic} & compositional & count or share & count & Ring-atom pairs between aromatic residues across TCR:peptide. \\
\texttt{ct\_\allowbreak{}tp\_\allowbreak{}hydrophobic} & compositional & count or share & count & Apolar C-C pairs between apolar residues across TCR:peptide. \\
\texttt{ct\_\allowbreak{}tp\_\allowbreak{}other} & compositional & count or share & count & Remaining classified TCR:peptide contacts. \\
\texttt{ct\_\allowbreak{}tm\_\allowbreak{}salt\_\allowbreak{}bridge} & compositional & count or share & count & Salt bridges across TCR:MHC. \\
\texttt{ct\_\allowbreak{}tm\_\allowbreak{}hydrogen\_\allowbreak{}bond} & compositional & count or share & count & Hydrogen bonds across TCR:MHC. \\
\texttt{ct\_\allowbreak{}tm\_\allowbreak{}aromatic} & compositional & count or share & count & Aromatic contacts across TCR:MHC. \\
\texttt{ct\_\allowbreak{}tm\_\allowbreak{}hydrophobic} & compositional & count or share & count & Hydrophobic contacts across TCR:MHC. \\
\texttt{ct\_\allowbreak{}tm\_\allowbreak{}other} & compositional & count or share & count & Remaining classified TCR:MHC contacts. \\
\texttt{cdr3\_\allowbreak{}dominance} & compositional & count or share & fraction & CDR3($\alpha$+$\beta$) share of all CDR TCR:peptide contacts. \\
\texttt{cdr3\_\allowbreak{}ab\_\allowbreak{}imbalance} & compositional & count or share & fraction & abs(CDR3a - CDR3b) / (CDR3a + CDR3b); how one-sided the CDR3 engagement is. \\
\texttt{chain\_\allowbreak{}cdr\_\allowbreak{}imbalance} & compositional & count or share & fraction & abs(a - b) / (a + b) over all CDR contacts; the chain-level mirror of chain\_balance. \\
\texttt{n\_\allowbreak{}clashes} & compositional & count or share & count & Peptide-partner heavy-atom pairs overlapping by more than 0.4 A on Bondi radii. \\
\texttt{clash\_\allowbreak{}score} & geometric & field moment & \AA & Summed overlap depth of those clashing pairs; the steric burden of a forced pose. \\
\texttt{mhc\_\allowbreak{}class\_\allowbreak{}bin}$^{\dagger}$ & categorical & count or share & class I / II & Which class of MHC presents the peptide: 0 for class I, 1 for class II. Class I and class II grooves differ in shape and in how they hold a peptide, so this conditions the other descriptors rather than being scored beside them; a coefficient fitted across both classes without it is fitted to a mixture. \\
\texttt{n\_\allowbreak{}loop\_\allowbreak{}contacts}$^{\dagger}$ & compositional & count or share & count & Contacts the six-CDR-loop partition sees; framework contacts are outside it by construction. \\
\texttt{n\_\allowbreak{}pep\_\allowbreak{}contacts} & compositional & count or share & count & Of those loop contacts, the ones reaching the peptide. \\
\texttt{n\_\allowbreak{}mhc\_\allowbreak{}contacts} & compositional & count or share & count & Of those loop contacts, the ones reaching the MHC. \\
\texttt{n\_\allowbreak{}pep\_\allowbreak{}int}$^{\dagger}$ & compositional & count or share & count & Intra-peptide residue contacts, at 5 A with a sequence separation of at least three. \\
\bottomrule
\end{longtable}

\subsection{The topology family (70 descriptors)}
\noindent The footprint: the contact graph, the cell tally, the two height fields, and the length-agnostic invariants of the C$\alpha$/C$\beta$ maps.

\small
\begin{longtable}{@{}>{\raggedright\arraybackslash}p{0.19\textwidth}p{0.14\textwidth}p{0.12\textwidth}p{0.08\textwidth}p{0.34\textwidth}@{}}
\toprule
descriptor & class & operator & units & definition \\
\midrule\endfirsthead
\toprule descriptor & class & operator & units & definition \\\midrule\endhead
\texttt{H\_\allowbreak{}cell} & compositional & Hill number & fraction & Normalized Shannon entropy of the contact composition over the twelve cells (six CDR loops x \{peptide, MHC\}). \\
\texttt{D1\_\allowbreak{}cell}$^{\dagger}$ & compositional & Hill number & count & Hill number of order 1 over the same cells, exp(H); monotone in H\_cell. \\
\texttt{D2\_\allowbreak{}cell} & compositional & Hill number & count & Hill number of order 2, 1/sum(p\textasciicircum{}2); the effective number of engaged cells, discounting the weakly populated ones. \\
\texttt{S\_\allowbreak{}cell} & compositional & Hill number & count & Richness: how many of the twelve cells are occupied. \\
\texttt{J\_\allowbreak{}cell}$^{\dagger}$ & compositional & Hill number & fraction & Pielou evenness over the occupied cells. NaN when one cell is occupied. \\
\texttt{H\_\allowbreak{}loop} & compositional & Hill number & fraction & Normalized entropy over the six CDR loops alone, ignoring which target each contact reaches. \\
\texttt{D2\_\allowbreak{}loop} & compositional & Hill number & count & Hill number of order 2 over the six loops. \\
\texttt{D2\_\allowbreak{}pep24} & compositional & Hill number & count & Hill number of order 2 over the twenty-four-cell partition, the peptide split into N-terminal, central and C-terminal bands. \\
\texttt{ab\_\allowbreak{}imb} & compositional & count or share & signed fraction & Signed (TRA - TRB)/(TRA + TRB) over CDR-loop contacts; positive is $\alpha$-shifted. \\
\texttt{ab\_\allowbreak{}imb\_\allowbreak{}pep} & compositional & count or share & signed fraction & The same restricted to peptide-side contacts. \\
\texttt{ab\_\allowbreak{}imb\_\allowbreak{}mhc} & compositional & count or share & signed fraction & The same restricted to MHC-side contacts. \\
\texttt{L\_\allowbreak{}canon} & compositional & count or share & log-odds & Canonical-docking log odds-ratio of loop class (germline, CDR3) against target (MHC, peptide), Haldane-Anscombe corrected. High when CDR3 sits on the peptide and the germline loops on the helices. \\
\texttt{p\_\allowbreak{}germ\_\allowbreak{}mhc} & compositional & count or share & fraction & Share of germline (CDR1/CDR2) contacts that reach the MHC. \\
\texttt{p\_\allowbreak{}cdr3\_\allowbreak{}pep} & compositional & count or share & fraction & Share of CDR3 contacts that reach the peptide. \\
\texttt{pep\_\allowbreak{}free\_\allowbreak{}frac} & compositional & count or share & fraction & Share of the peptide the groove leaves for the receptor: mean over positions of n\_TCR/(n\_TCR + n\_MHC). The threshold-free reading of 'peptide without its MHC anchors'. \\
\texttt{pep\_\allowbreak{}cov\_\allowbreak{}frac} & compositional & count or share & fraction & Peptide positions the TCR contacts, over peptide length. \\
\texttt{pep\_\allowbreak{}cov\_\allowbreak{}even} & compositional & Hill number & fraction & Pielou evenness of the accessibility-discounted contact distribution, base ln(peptide length); how evenly the receptor uses the peptide it can reach. \\
\texttt{pep\_\allowbreak{}cov\_\allowbreak{}d2n} & compositional & Hill number & fraction & Hill number of order 2 of that distribution over peptide length; the effective share of the peptide engaged. \\
\texttt{pep\_\allowbreak{}cov\_\allowbreak{}centre} & compositional & count or share & fraction & Contact-weighted mean position on [0, 1] from N- to C-terminus; 0.5 is centred. \\
\texttt{pep\_\allowbreak{}cov\_\allowbreak{}spread} & compositional & count or share & fraction & Contact-weighted standard deviation of that position, doubled; approaches 1 when the receptor reaches both termini. \\
\texttt{h0\_\allowbreak{}pers\_\allowbreak{}ent} & geometric & Hill number & fraction & Normalized entropy of the H0 barcode of the contacted pMHC C$\alpha$ cloud. The bar lengths are the minimum spanning tree's edges, so no filtration is chosen. \\
\texttt{g\_\allowbreak{}even\_\allowbreak{}tcr}$^{\dagger}$ & topological & Hill number & fraction & Pielou evenness of the contact degrees of the engaged CDR-loop residues, base the engaged count. 1 when every engaged residue carries the same number of partners. \\
\texttt{g\_\allowbreak{}even\_\allowbreak{}pmhc} & topological & Hill number & fraction & Pielou evenness of the contact degrees of the engaged pMHC residues, base the engaged count. \\
\texttt{g\_\allowbreak{}comp\_\allowbreak{}frac}$^{\dagger}$ & topological & homology invariant & fraction & Connected components of the bipartite contact graph per node; the parameter-free form of the footprint patch count, needing no C$\alpha$ radius. \\
\texttt{g\_\allowbreak{}alg\_\allowbreak{}conn} & topological & spectral invariant & ratio & Second-smallest eigenvalue of the normalised Laplacian on the largest contact-graph component, in [0, 2]. Near 0 when the footprint is about to fall into two patches. \\
\texttt{g\_\allowbreak{}cyclo\_\allowbreak{}frac} & topological & homology invariant & fraction & Contacts beyond a spanning forest over all contacts, (E - V + C) / E, which is the contact graph's first Betti number made size-free. High when the footprint is interlocked. \\
\texttt{g\_\allowbreak{}loop\_\allowbreak{}even}$^{\dagger}$ & compositional & Hill number & fraction & Pielou evenness over the six CDR loops of the number of distinct pMHC residues each reaches, base 6. Counts partners rather than contacts, so it does not track residue size. \\
\texttt{g\_\allowbreak{}loop\_\allowbreak{}overlap} & compositional & count or share & fraction & Mean pairwise Jaccard overlap of the engaged CDR loops' pMHC partner sets. High when the loops crowd onto the same residues instead of partitioning the surface. \\
\texttt{g\_\allowbreak{}assort} & topological & correlation & ratio & Degree assortativity of the contact graph: the correlation, over contacts, between the degrees of the two residues involved. \\
\texttt{degree\_\allowbreak{}evenness\_\allowbreak{}tp}$^{\dagger}$ & compositional & Hill number & fraction & Participation ratio of the receptor-side contact degrees across TCR:peptide, in [0, 1]. Low when a few over-reaching side chains hoard the contact budget. \\
\texttt{frac\_\allowbreak{}well\_\allowbreak{}coordinated\_\allowbreak{}tp}$^{\dagger}$ & compositional & count or share & fraction & Share of contacting receptor residues reaching no more than three peptide residues, the count a crystal side chain typically makes. \\
\texttt{m\_\allowbreak{}erank\_\allowbreak{}tp}$^{\dagger}$ & geometric & Hill number & fraction & Effective rank of the CDR-loop x peptide C$\alpha$ proximity kernel over its maximum: how many independent approach modes the interface has. Peptide-length coupled (Spearman -0.547 on 148 class I crystals) because a longer class I peptide bulges. \\
\texttt{m\_\allowbreak{}gap\_\allowbreak{}tp}$^{\dagger}$ & geometric & spectral invariant & ratio & Second over first singular value of that kernel. Near 0 when the approach is separable into a loop profile times a peptide profile rather than pairing specific residues. \\
\texttt{m\_\allowbreak{}erank\_\allowbreak{}tm}$^{\dagger}$ & geometric & Hill number & fraction & Effective rank fraction of the CDR-loop x MHC-helix kernel. CDR3-length coupled (Spearman -0.437), so it reads how much loop there is to spread over the helices. \\
\texttt{m\_\allowbreak{}gap\_\allowbreak{}tm} & geometric & spectral invariant & ratio & Second over first singular value of the CDR-loop x MHC-helix kernel. The least length-coupled column in the catalogue: -0.012 against CDR3 length, +0.027 against peptide length, 99.9 per cent of its variance surviving both. \\
\texttt{m\_\allowbreak{}face\_\allowbreak{}tp} & geometric & field moment & \AA & Mean C$\alpha$-C$\alpha$ minus C$\beta$-C$\beta$ distance over the contacting TCR:peptide residue pairs. Positive when side chains lean towards each other, negative when the backbones are the close part and the side chains point away. \\
\texttt{m\_\allowbreak{}face\_\allowbreak{}tm} & geometric & field moment & \AA & Mean C$\alpha$-C$\alpha$ minus C$\beta$-C$\beta$ distance over the contacting TCR:MHC residue pairs. \\
\texttt{ca\_\allowbreak{}cb\_\allowbreak{}agreement\_\allowbreak{}tp} & geometric & correlation & ratio & Spearman correlation between the C$\alpha$ and C$\beta$ distance maps over the TCR:peptide approach shell. High when the side chains track the backbone, as they do in a crystal. \\
\texttt{ca\_\allowbreak{}cb\_\allowbreak{}agreement\_\allowbreak{}tm}$^{\dagger}$ & geometric & correlation & ratio & The same rank correlation across the TCR:MHC approach shell. \\
\texttt{fp\_\allowbreak{}b0\_\allowbreak{}r7} & topological & homology invariant & count & Betti-0 of the flag complex on the contacted pMHC C$\alpha$ atoms at 7 A: how many disconnected patches the footprint falls into. \\
\texttt{fp\_\allowbreak{}b1\_\allowbreak{}r7} & topological & homology invariant & count & Betti-1 at 7 A: how many holes the footprint encloses. \\
\texttt{fp\_\allowbreak{}chi\_\allowbreak{}r7}$^{\dagger}$ & topological & homology invariant & count & Euler characteristic b0 - b1 at 7 A. \\
\texttt{fp\_\allowbreak{}b0\_\allowbreak{}frac\_\allowbreak{}r7} & topological & homology invariant & fraction & Patches per contacted residue at 7 A; the size-free form of Betti-0. \\
\texttt{fp\_\allowbreak{}b0\_\allowbreak{}r8} & topological & homology invariant & count & Betti-0 of the flag complex on the contacted pMHC C$\alpha$ atoms at 8 A: how many disconnected patches the footprint falls into. \\
\texttt{fp\_\allowbreak{}b1\_\allowbreak{}r8} & topological & homology invariant & count & Betti-1 at 8 A: how many holes the footprint encloses. \\
\texttt{fp\_\allowbreak{}chi\_\allowbreak{}r8}$^{\dagger}$ & topological & homology invariant & count & Euler characteristic b0 - b1 at 8 A. \\
\texttt{fp\_\allowbreak{}b0\_\allowbreak{}frac\_\allowbreak{}r8} & topological & homology invariant & fraction & Patches per contacted residue at 8 A; the size-free form of Betti-0. \\
\texttt{sc\_\allowbreak{}shape} & geometric & correlation & ratio & Pearson r between the pMHC and TCR height fields over the shared grid; positive is complementary, the receptor riding up where the groove rises. Lawrence \& Colman's Sc is the same idea on a dot surface. \\
\texttt{sc\_\allowbreak{}charge} & compositional & correlation & ratio & Pearson r between the two charge fields; NEGATIVE is complementary, plus meeting minus. \\
\texttt{sc\_\allowbreak{}phobic} & compositional & correlation & ratio & Pearson r between the two Kyte-Doolittle fields; positive is complementary, apolar meeting apolar. \\
\texttt{sc\_\allowbreak{}charge\_\allowbreak{}prod} & compositional & field moment & ratio & Mean per-cell product of the two charge fields. \\
\texttt{sc\_\allowbreak{}phobic\_\allowbreak{}prod} & compositional & field moment & ratio & Mean per-cell product of the two hydropathy fields. \\
\texttt{sc\_\allowbreak{}gap\_\allowbreak{}mean} & geometric & field moment & \AA & Mean of h(TCR) - h(pMHC) over retained cells. Negative on a real interface: the median cell interdigitates. \\
\texttt{sc\_\allowbreak{}gap\_\allowbreak{}sd} & geometric & field moment & \AA & Spread of the same gap. High when the receptor rests on a few high points rather than meshing. \\
\texttt{sc\_\allowbreak{}gap\_\allowbreak{}vol} & geometric & field moment & \AA$^3$ & Void volume, the gap integrated over the contact plane where it is positive. \\
\texttt{sc\_\allowbreak{}interlock} & geometric & field moment & \AA$^3$ & Interdigitated volume, the gap integrated where it is negative. The larger of the two on a real interface. \\
\texttt{sc\_\allowbreak{}gap\_\allowbreak{}index} & geometric & field moment & \AA & Void volume over retained contact area; the intensive form of the gap-volume channel. \\
\texttt{sc\_\allowbreak{}interlock\_\allowbreak{}frac} & geometric & field moment & fraction & Share of retained cells whose gap is negative; the per-structure form of the corpus 71\% interdigitation. \\
\texttt{sc\_\allowbreak{}gap\_\allowbreak{}depth} & geometric & field moment & \AA & Mean depth over the interlocked cells alone: how far the receptor reaches in where it does. \\
\texttt{sc\_\allowbreak{}gap\_\allowbreak{}height} & geometric & field moment & \AA & Mean standoff over the void cells alone: how high it stands where it does not mesh. \\
\texttt{sc\_\allowbreak{}gap\_\allowbreak{}asym} & geometric & field moment & signed fraction & (void - interlock) / (void + interlock); -1 for a face that only interlocks, +1 for one that only stands off. \\
\texttt{sc\_\allowbreak{}dh} & geometric & field moment & \AA & Mean absolute per-cell height difference between the two faces. \\
\texttt{sc\_\allowbreak{}dcharge} & compositional & field moment & ratio & Mean absolute per-cell charge difference between the two faces. \\
\texttt{sc\_\allowbreak{}dphobic} & compositional & field moment & ratio & Mean absolute per-cell hydropathy difference between the two faces. \\
\texttt{sc\_\allowbreak{}cells} & geometric & count or share & count & Grid cells entering the comparison; bookkeeping, so a low complementarity can be told from a thin one. \\
\texttt{sc\_\allowbreak{}coverage} & geometric & count or share & fraction & Retained cells as a share of the occupied pMHC cells in the window; bookkeeping. \\
\texttt{co\_\allowbreak{}pep} & compositional & count or share & ratio & Contact order on the peptide: mean sequence separation of the peptide residues one CDR loop reaches, averaged over loops and divided by the peptide's span. \\
\texttt{co\_\allowbreak{}mhc} & compositional & count or share & ratio & Contact order on the MHC helices, by the same construction. \\
\texttt{partcoef\_\allowbreak{}tcr} & compositional & Hill number & fraction & Mean over engaged TCR residues of 1 - sum\_s (k\_s/k)\textasciicircum{}2 with the modules peptide and MHC; 0 when every residue reads one target only. \\
\texttt{partcoef\_\allowbreak{}pmhc} & compositional & Hill number & fraction & The same over engaged pMHC residues with the six CDR loops as modules. \\
\bottomrule
\end{longtable}

\subsection{The energetics family (15 descriptors)}
\noindent Sums of a pair potential over a contact set, and differences of them.

\small
\begin{longtable}{@{}>{\raggedright\arraybackslash}p{0.19\textwidth}p{0.14\textwidth}p{0.12\textwidth}p{0.08\textwidth}p{0.34\textwidth}@{}}
\toprule
descriptor & class & operator & units & definition \\
\midrule\endfirsthead
\toprule descriptor & class & operator & units & definition \\\midrule\endhead
\texttt{Phi\_\allowbreak{}tcr\_\allowbreak{}pep} & energetic & potential sum & log-odds & $\Phi$ over TCR-peptide contacts under TCRen2, summed over all TCR regions. Lower is more favourable. \\
\texttt{Phi\_\allowbreak{}tcr\_\allowbreak{}mhc} & energetic & potential sum & log-odds & $\Phi$ over TCR-MHC contacts under Miyazawa-Jernigan. \\
\texttt{Phi\_\allowbreak{}cdr12} & energetic & potential sum & log-odds & The CDR1 + CDR2 part of the TCR:peptide energy, both chains. \\
\texttt{Phi\_\allowbreak{}cdr3a} & energetic & potential sum & log-odds & The CDR3$\alpha$ part of the TCR:peptide energy. \\
\texttt{Phi\_\allowbreak{}cdr3b} & energetic & potential sum & log-odds & The CDR3$\beta$ part of the TCR:peptide energy. \\
\texttt{dPhi\_\allowbreak{}tcr\_\allowbreak{}pep} & energetic & potential sum & log-odds & Poly-alanine reference delta of the TCR:peptide energy; the pose-geometry baseline removed. \\
\texttt{dPhi\_\allowbreak{}pep\_\allowbreak{}soft} & energetic & potential sum & log-odds & Smoothed reference delta, peptide direction: the peptide's energy minus the free energy of the residue background at each peptide position, receptor frozen. \\
\texttt{varPhi\_\allowbreak{}pep\_\allowbreak{}soft} & energetic & potential sum & log-odds$^2$ & Variance of the local field under the background, peptide direction: how sharply each peptide position's energy responds to residue identity, summed over positions. NOT a ddG -- it is a second cumulant, not a difference of differences. \\
\texttt{dPhi\_\allowbreak{}tcr\_\allowbreak{}soft} & energetic & potential sum & log-odds & Smoothed reference delta, receptor direction: the receptor's energy minus the free energy of the residue background at each contacted TCR position, peptide frozen. \\
\texttt{varPhi\_\allowbreak{}tcr\_\allowbreak{}soft} & energetic & potential sum & log-odds$^2$ & Variance of the local field under the background, receptor direction, summed over contacted TCR positions. \\
\texttt{dPhi\_\allowbreak{}tra\_\allowbreak{}soft} & energetic & potential sum & log-odds & The $\alpha$-chain part of the receptor-direction smoothed reference delta. \\
\texttt{dPhi\_\allowbreak{}trb\_\allowbreak{}soft} & energetic & potential sum & log-odds & The $\beta$-chain part of the receptor-direction smoothed reference delta. \\
\texttt{Phi\_\allowbreak{}pep\_\allowbreak{}mhc}$^{\dagger}$ & energetic & potential sum & log-odds & $\Phi$ over peptide-MHC contacts under Miyazawa-Jernigan. Computed without the receptor. \\
\texttt{dPhi\_\allowbreak{}pep\_\allowbreak{}mhc}$^{\dagger}$ & energetic & potential sum & log-odds & The same reference across peptide:MHC. Computed without the receptor. \\
\texttt{Phi\_\allowbreak{}pep\_\allowbreak{}int}$^{\dagger}$ & energetic & potential sum & log-odds & The peptide's own intra-chain contact energy. Computed without the receptor. \\
\bottomrule
\end{longtable}

\subsection{The potts family (5 descriptors)}
\noindent The same interface read against the partition function of the coupled contact model.

\small
\begin{longtable}{@{}>{\raggedright\arraybackslash}p{0.19\textwidth}p{0.14\textwidth}p{0.12\textwidth}p{0.08\textwidth}p{0.34\textwidth}@{}}
\toprule
descriptor & class & operator & units & definition \\
\midrule\endfirsthead
\toprule descriptor & class & operator & units & definition \\\midrule\endhead
\texttt{neg\_\allowbreak{}energy}$^{\dagger}$ & energetic & potential sum & kT & -E of the observed contact map under the coupled Potts model; higher is more native-like. Exactly log\_z + log\_lik. \\
\texttt{log\_\allowbreak{}z} & energetic & potential sum & kT & Log partition function over every contact map the geometry admits; the interface's capacity. \\
\texttt{log\_\allowbreak{}lik} & energetic & potential sum & kT & Log probability of the observed contact map; its typicality. \\
\texttt{psi} & energetic & potential sum & kT/site & log\_lik per available site, so interfaces of different size compare. \\
\texttt{n\_\allowbreak{}contacts} & energetic & count or share & count & Available residue pairs that engaged. Distinct from the footprint's loop tally. \\
\bottomrule
\end{longtable}

\subsection{The kinetics family (17 descriptors)}
\noindent The contact map as a network of breakable springs. No potential enters.

\small
\begin{longtable}{@{}>{\raggedright\arraybackslash}p{0.19\textwidth}p{0.14\textwidth}p{0.12\textwidth}p{0.08\textwidth}p{0.34\textwidth}@{}}
\toprule
descriptor & class & operator & units & definition \\
\midrule\endfirsthead
\toprule descriptor & class & operator & units & definition \\\midrule\endhead
\texttt{exp\_\allowbreak{}lost} & geometric & field moment & count & Expected TCR:peptide contacts lost under a 1 A isotropic shift. \\
\texttt{mean\_\allowbreak{}margin} & geometric & field moment & \AA & Mean contact margin, cutoff minus minimum heavy-atom distance. \\
\texttt{frac\_\allowbreak{}robust} & compositional & count or share & fraction & Share of TCR:peptide contacts with at least 1 A of margin. \\
\texttt{K\_\allowbreak{}tens} & geometric & spectral invariant & N/m & Tensile stiffness along the docking axis. \\
\texttt{K\_\allowbreak{}shear} & geometric & spectral invariant & N/m & In-plane stiffness, S\_tot minus K\_tens. \\
\texttt{aniso}$^{\dagger}$ & geometric & spectral invariant & ratio & K\_shear / K\_tens; how much stiffer the interface is along the pull than across it. \\
\texttt{lam\_\allowbreak{}max} & geometric & spectral invariant & N/m & Largest eigenvalue of the stiffness tensor. \\
\texttt{lam\_\allowbreak{}min} & geometric & spectral invariant & N/m & Smallest eigenvalue of the stiffness tensor. \\
\texttt{n\_\allowbreak{}spring} & compositional & count or share & count & Springs in the interface network; every other kinetics column is NaN below three. \\
\texttt{S\_\allowbreak{}tot}$^{\dagger}$ & geometric & spectral invariant & N/m & Trace of the stiffness tensor; total interface stiffness. \\
\texttt{rupture\_\allowbreak{}force} & geometric & mechanical work & N & Peak resisting force under steered separation along the weaker axis. \\
\texttt{rupture\_\allowbreak{}work} & geometric & mechanical work & J & Force integrated to full separation; the off-rate proxy, and a geometry-only quantity no potential enters. \\
\texttt{couple\_\allowbreak{}pep} & compositional & count or share & count & Peptide residues contacting both the MHC and the TCR. \\
\texttt{couple\_\allowbreak{}mhc} & compositional & count or share & count & MHC residues contacting both the peptide and the TCR. \\
\texttt{couple\_\allowbreak{}tcr} & compositional & count or share & count & TCR residues in the V$\alpha$-V$\beta$ interface that also contact the pMHC. \\
\texttt{couple\_\allowbreak{}total}$^{\dagger}$ & compositional & count or share & count & Sum of the three coupling counts. \\
\texttt{n\_\allowbreak{}interface} & compositional & count or share & count & Interface residue count; the size denominator for the coupling counts. \\
\bottomrule
\end{longtable}

\noindent $^{\dagger}$ carries a \texttt{STATUS} flag; see \cref{sec:status}.
